\documentclass[10pt]{article}
\usepackage[margin=1.15in]{geometry}
\usepackage[T1]{fontenc}
\usepackage{lmodern,microtype}
\usepackage{amsmath,amssymb,amsthm,mathtools}
\usepackage[hidelinks]{hyperref}

\setlength{\parindent}{0pt}
\setlength{\parskip}{0.34em}
\setlength{\abovedisplayskip}{5pt plus 1pt minus 2pt}
\setlength{\belowdisplayskip}{5pt plus 1pt minus 2pt}
\setlength{\abovedisplayshortskip}{3pt plus 1pt}
\setlength{\belowdisplayshortskip}{3pt plus 1pt}

\newtheoremstyle{compact}
  {5pt}{5pt}{\itshape}{}{}{.}{0.45em}{}
\theoremstyle{compact}
\newtheorem{theorem}{Theorem}

\newcommand{\EP}{N}
\newcommand{\C}{\mathfrak C}

\begin{document}

\begin{center}
{\LARGE\bfseries Finite-Source Causal Dilution}\par
\vspace{0.18em}
{\large Repeated training cannot preserve information concentration}\par
\vspace{0.55em}
{Jonathan R. Landers \qquad July 2026}
\end{center}
\vspace{0.25em}

A model may circle through the same dataset many times. Each lap changes the weights, lowers loss, and may sharpen recall, so it is natural to imagine that every epoch contributes a fresh packet of information. But the source has not changed. Repetition can move information from the dataset into the training trajectory; it cannot create more information about that dataset than the dataset contains.

There is a second law. A finite response is not instantaneous: the effect of a training round is distributed over later steps. If successive rounds act as passive causal stages, their response kernels compose by convolution and spread strictly in entropy power. The theorem below is the collision of these two facts: \emph{the information reservoir is finite, while the causal response width keeps growing.}

\paragraph{Training information.}
Let $D$ be a random finite dataset and $W_0,W_1,\ldots,W_n$ the parameter states generated by repeatedly training on $D$. Relative to the initial state, define
\[
J_n:=I(D;W_{1:n}\mid W_0),
\qquad W_{1:n}:=(W_1,\ldots,W_n).
\]
The chain rule gives
\[
J_n=\sum_{t=1}^{n}I(D;W_t\mid W_{0:t-1})
\le H(D\mid W_0).
\tag{1}
\]
Thus many rounds may contribute, but the cumulative contribution draws from one finite conditional-information budget.

\paragraph{Finite response.}
Let $K_t$ be a nondegenerate probability density on $[0,\infty)$ describing the lagged response of round $t$. Independent sequential stages compose as
\[
K_{1:n}:=K_n*K_{n-1}*\cdots*K_1.
\tag{2}
\]
For a density $K$ with differential entropy $h(K)$, define its entropy power
\[
\EP(K):=\frac{1}{2\pi e}\exp\{2h(K)\}.
\tag{3}
\]
All logarithms are natural. The quantity $\sqrt{\EP(K)}$ is an entropic response-time scale.

\begin{theorem}[Finite-source causal dilution]
Assume $H(D\mid W_0)<\infty$. Let $K_1,\ldots,K_n$, $n\ge2$, be independent-stage passive causal densities with finite differential entropy, and define $K_{1:n}$ by \emph{(2)}. Set
\[
\C_n:=\frac{I(D;W_{1:n}\mid W_0)}{\sqrt{\EP(K_{1:n})}}.
\]
Then
\[
\boxed{
\C_n<
\frac{H(D\mid W_0)}
{\sqrt{\displaystyle\sum_{t=1}^{n}\EP(K_t)}}.}
\tag{4}
\]
Consequently, if $\EP(K_t)\ge\nu>0$ for every round, then
\[
\boxed{
\C_n<\frac{H(D\mid W_0)}{\sqrt{n\nu}}
\longrightarrow0.}
\tag{5}
\]
\end{theorem}

\newpage
\section*{Proof}
The numerator and denominator obey different laws.

For the numerator,
\[
I(D;W_{1:n}\mid W_0)
=H(D\mid W_0)-H(D\mid W_{0:n})
\le H(D\mid W_0).
\tag{6}
\]
This identity is the fixed-source budget. It permits gradual extraction, compression, and memorization; it forbids only the manufacture of new information about $D$ by exposing the learner to the same $D$ again.

For the denominator, let $T_t\ge0$ be independent random delays with densities $K_t$. Their sum $T_{1:n}=T_1+\cdots+T_n$ has density $K_{1:n}$. The Shannon--Stam entropy-power inequality yields
\[
\EP(K_{1:n})\ge\sum_{t=1}^{n}\EP(K_t).
\tag{7}
\]
The inequality is strict. Equality would require Gaussian summands, but a nondegenerate Gaussian has support on all of $\mathbb R$, whereas each passive causal stage is supported on $[0,\infty)$. Hence
\[
\EP(K_{1:n})>\sum_{t=1}^{n}\EP(K_t).
\tag{8}
\]
Combining \emph{(6)} and \emph{(8)},
\[
\C_n
=\frac{I(D;W_{1:n}\mid W_0)}{\sqrt{\EP(K_{1:n})}}
<\frac{H(D\mid W_0)}
{\sqrt{\sum_{t=1}^{n}\EP(K_t)}},
\]
which proves \emph{(4)}. If $\EP(K_t)\ge\nu$, then the denominator exceeds $\sqrt{n\nu}$, proving \emph{(5)}. \hfill$\square$

\section*{Interpretation}
The result does \emph{not} say that later epochs are useless, nor that information in the weights must fall. Indeed, $J_n$ can increase as the trajectory captures more of what was already present in $D$. It says something more geometric:
\[
\boxed{\text{extraction is source-limited, while finite causal mediation is spread-producing.}}
\]
The numerator may approach a ceiling. The denominator must grow whenever every round contributes a genuinely finite passive response stage. Therefore no such loop can sustain a nonzero amount of acquired dataset information per unit entropy-power response width indefinitely.

This gives a concrete prediction for checkpointed language-model training. Estimate $J_n$ across perturbed datasets or ensembles of runs, and estimate $K_{1:n}$ from the lagged influence of examples on losses, logits, or representations. The signature is a separation of scales:
\[
J_n\ \text{levels off or grows slowly},
\qquad
\EP(K_{1:n})\ \text{continues to grow}.
\]
Late training is then not information arriving from another lap around the same source. It is the redistribution, reinforcement, compression, or memorization of a finite source through an increasingly extended causal response.

\vfill
{\footnotesize\textbf{Context.} The kernel law is the passive-causal composition framework of J.~R.~Landers, \emph{The Entropy of Finite Response} (2026). The proof synthesizes its strict causal entropy-power spreading with the conditional mutual-information budget of a fixed dataset.}

\end{document}
